\section{Experimental details}\label{app:details}

\subsection{Rollout invariance (\texorpdfstring{\cref{sec:s1}}{Section 6.1})}
\begin{itemize}[leftmargin=*,itemsep=2pt]
\item \textbf{Checkpoints.}
  \begin{itemize}[leftmargin=*,itemsep=1pt]
  \item The UAQ checkpoints divide every $W_{Q,K,V}$ by $s$ and multiply every attention input-norm gain by $s$. Both checkpoints are stored in BF16, like the original.
  \item At $s=2$, no INT8 or FP8 code changes, and the stored INT8 scales are exactly halved.
  \item At $s=1.5$, BF16 storage re-rounds 4.0--4.3\% of the INT8 per-channel codes and 2.0--2.6\% of the FP8 per-tensor codes. Quantized checkpoints whose scales are also stored in BF16 change 5.4--5.7\% of their codes.
  \end{itemize}
\item \textbf{Serving.} vLLM 0.33.0. The logs confirm \texttt{CutlassInt8ScaledMMLinearKernel} for INT8 W8A8 and \texttt{CutlassFP8ScaledMMLinearKernel} for both FP8 modes.
\item \textbf{Responses.} 512 GSM8K training prompts, one response each, sampled from the original BF16 model at temperature 1. Every response token is teacher-forced under every configuration.
\item \textbf{Metrics.}
  \begin{itemize}[leftmargin=*,itemsep=1pt]
  \item The primary metric is the mean $|\log\mu-\log\pi|$ over sampled tokens, with $\pi$ the BF16 policy of the same checkpoint.
  \item The KL estimators are importance-weighted when the reference is the UAQ checkpoint's BF16 policy, because the tokens were sampled from the original one.
  \item Intervals: paired bootstrap over the 512 responses.
  \end{itemize}
\item \textbf{Why mean $|\Delta\log p|$ is primary.} Both sampled-token KL estimators rely on $\E_\pi[r-1]=0$. A few rare tokens with $\mu\gg\pi$ (with $r$ up to 17--37) make the sample mean of $r-1$ unreliable for FP8 per-tensor at 1.5B. Removing the 100 largest values of $r-1$ (0.05\% of tokens) changes it from $+1.7\times10^{-3}$ to $-2.8\times10^{-4}$.
\item \textbf{Determinism.} Every configuration was re-run with shuffled batches, and the FP8 runs at 0.5B and 1.5B were repeated in the same batch order. At 7B, every configuration is batch-invariant. Each model took 14--18 minutes on the H100.
\end{itemize}

\subsection{Exact equivalence in a training loop (\texorpdfstring{\cref{sec:s2}}{Section 6.2})}
\begin{itemize}[leftmargin=*,itemsep=2pt]
\item \textbf{Loop.} PyTorch 2.14 on the H100. GSM8K training prompts; 20 rollout batches of 8 prompts $\times$ 4 samples, with at most 256 new tokens.
  \begin{itemize}[leftmargin=*,itemsep=1pt]
  \item GRPO advantages, and the decoupled PPO loss with clip 0.2, TIS cap $C=2$ and the ACR bound.
  \item Two epochs of two minibatches per rollout batch, so 4 AdamW updates per batch and 80 in all.
  \item $\lambda=0.01$, $\varepsilon=10^{-8}$, global-norm clipping at $\tau=0.1$, seed 0.
  \end{itemize}
\item \textbf{Rollouts.} INT8 W8A8 fake quantization in PyTorch, with per-channel absmax weights and dynamic per-token absmax activations. Tokens are sampled by inverse CDF from seeded uniforms. vLLM is not used, because it re-quantizes only FP8 online and its default engine is not repeatable at this size (\cref{sec:s1}).
\item \textbf{Training numerics.} BF16 autocast with FP32 master weights; PyTorch's deterministic-algorithms mode; eager attention; TF32 off.
\item \textbf{Fingerprints.} At every step, SHA-256 hashes of the sampled tokens, $\mu$, $\pi_{\mathrm{old}}$, and the per-group weights mapped to original coordinates. Weights were also saved after steps 1 and 20.
\item \textbf{What was exercised} (R4).
  \begin{itemize}[leftmargin=*,itemsep=1pt]
  \item Clipping bound on every update with a nonzero gradient (gradient norms 0.27--3.2 against $\tau=0.1$).
  \item The remaining updates had an exactly zero gradient, because every group in the minibatch had uniform rewards; AdamW still moved the weights through its momentum.
  \item TIS truncated 0.2--1.3\% of tokens per batch, and ACR's raised bound kept 26 ($\eta=10^{-6}$) and 7 ($\eta=10^{-5}$) token terms that PPO clipping would have dropped.
  \item Each arm took 7.7 minutes.
  \end{itemize}
\end{itemize}

\subsection{Implementation in VeRL}\label{app:verl}
\begin{itemize}[leftmargin=*,itemsep=2pt]
\item \textbf{Stack.} VeRL v0.9.1 with its locked dependencies: vLLM 0.24.0, PyTorch 2.11 and Transformers 5.9.0. Actor and critic use PyTorch SDPA attention without sequence packing, because the pinned flash-attention wheel is no longer available. A two-line patch lets VeRL's critic loader honour the configured attention implementation; it is our only change to VeRL's code.
\item \textbf{INT8 rollouts re-quantized at every sync.} VeRL natively refits only FP8 blockwise rollouts.
  \begin{itemize}[leftmargin=*,itemsep=1pt]
  \item vLLM is started with a compressed-tensors INT8 W8A8 configuration: per-channel weights, dynamic per-token activations, and the output head left in BF16.
  \item At every sync, a vLLM worker extension quantizes each incoming linear weight row by row: scale $=$ row absmax $\times$ fl32$(1/127)$, and codes $=\mathrm{round}(W/\text{scale})$ clamped to $\pm127$. It writes codes and scales into the live INT8 parameters in the layout of the CUTLASS kernel.
  \item The synced weights arrive as the FP32 master weights, so quantization starts from FP32 without BF16 rounding first.
  \item The scale is computed as a multiplication on purpose: PyTorch divides a GPU tensor by a scalar through its reciprocal, which differs by one ulp from true division in about 5\% of rows.
  \item vLLM's post-load weight processing is disabled during syncs.
  \item We verified that all 168 linear weights of the 0.5B model are re-quantized at every sync, and that the resulting codes and scales are bit-identical to an independent quantization.
  \end{itemize}
\item \textbf{ACR} is a custom policy loss. With no truncated token, it is bit-identical to VeRL's PPO loss with TIS weights. With truncation, it matches a float64 reference of QuRL's ACR loss to $6\times10^{-10}$.
\item \textbf{Matched settings (R5).} An optimizer plug-in for VeRL's FSDP2 actor builds three parameter groups: Q/K/V weights, attention input-norm gains, and everything else.
  \begin{itemize}[leftmargin=*,itemsep=1pt]
  \item Each group gets its own learning rate, $\varepsilon$ and weight decay from \cref{thm:reparam}(b).
  \item It replaces the clipping step with rule (c): the clipping norm is computed on $(s\,g_W,g_\gamma/s,g_{\mathrm{rest}})$ for R5, and on the ordinary gradient for every other arm.
  \item So all six arms run identical optimizer and clipping code.
  \item VeRL's constant learning-rate scheduler keeps each group's rate.
  \end{itemize}
\item \textbf{Exactness check in VeRL.} The plug-in was checked on the H100 inside VeRL's FSDP2 engine (\cref{sec:s2}): R5 and an R4 repeat were bit-identical to R4 after three updates. Each run of the UAQ experiment starts only after this check passes.
\end{itemize}

\subsection{GSM8K PPO configuration (\texorpdfstring{\cref{sec:s3}}{Section 6.3})}
\begin{center}\small
\begin{tabularx}{\linewidth}{@{}l L@{}}
\toprule
Setting & Value\\
\midrule
Model & Qwen2.5-0.5B-Instruct (actor, reference and critic initialization)\\
Data & GSM8K: 7{,}473 training and 1{,}319 test problems; VeRL's prompt template and strict answer format\\
Algorithm & PPO with a critic; GAE with $\gamma=\lambda=1$; one PPO epoch; clip 0.2; dual clip 3.0; token-mean loss\\
Batch & 256 prompts per step, PPO mini-batch 64 (4 updates per step), 15 epochs = 435 steps\\
Lengths & prompts up to 320 tokens (all GSM8K prompts fit in 256), responses up to 512\\
Actor optimizer & AdamW, betas (0.9, 0.999), weight decay 0.01, gradient clip 1.0, constant learning rate without warmup: $10^{-6}$ (UAQ experiment) or $10^{-5}$ (reproduction)\\
Critic & AdamW, learning rate $10^{-5}$, value clip 0.5\\
Regularizers & KL penalty to the reference in the reward, 0.001 ($k_1$, fixed); entropy bonus 0.001\\
Rollouts & vLLM 0.24, temperature 1; validation greedy every 29 steps, with the rollout engine (INT8 for the QuRL arms)\\
QuRL arms & INT8 W8A8 re-quantized at every sync; token-level TIS, $C=2$; ACR\\
UAQ & $s=1.5$ on Q/K/V and the attention input-norm gains\\
Seeds & seed $k$ sets the data order and the rollout seed $42+k$; the arms of a seed share both\\
Run time & 4.2--5.5\,h per run in the UAQ experiment, 104\,h for its 22 runs; 2.9\,h and 4.9\,h for the two reproduction runs\\
\bottomrule
\end{tabularx}
\end{center}

\subsection{Tracking and persistence run (\texorpdfstring{\cref{sec:e2}}{Section 6.4})}\label{app:e2-details}
\begin{itemize}[leftmargin=*,itemsep=2pt]
\item \textbf{Recording.} A subclass of the INT8 refit extension (\cref{app:verl}), otherwise unchanged, copies each synced FP32 linear weight before quantizing it. Sync $k$ carries the weights after step $k-1$.
  \begin{itemize}[leftmargin=*,itemsep=1pt]
  \item Full states, about 2\,GB each, at the 14 syncs of \cref{sec:e2}.
  \item At every sync, 16 fixed rows of each of the seven linear weights of layers 0, 6, 12, 18 and 23: 818{,}630 weights once the absmax elements are removed.
  \item At each full sync, layer~0's live Q/K/V codes and scales.
  \end{itemize}
  The records take 35\,GB; the run took 4.5\,h.
\item \textbf{Dump check.} At each of the 14 full syncs, every linear weight is present in FP32 under the same name as at sync~1. Layer~0's live INT8 codes and scales equal the quantizer applied to the record, and the row samples equal the matching rows of the full record.
\item \textbf{Ties.} 0.18\% of the weights of the BF16 checkpoint lie exactly on a rounding boundary ($|\varphi|=\tfrac12$), against $2\times10^{-6}$ after one step (FP32 round-off). An update moves a tie to one side or the other, so it flips with probability about $\tfrac12$ whatever its size. The registered rules for the first pairs therefore leave out the weights that were ties at sync~1.
\item \textbf{Capture statistics,} for each one-step pair:
  \begin{itemize}[leftmargin=*,itemsep=1pt]
  \item the flip share $F$ against $P=\mathrm{mean}\min(1,|u|)$;
  \item the calibration slope, a count-weighted slope through the origin of observed against predicted flip rates in quarter-decade bins of $|u|$ with at least $10^4$ weights;
  \item the capture slope $\sum\Delta c\,u/\sum u^2$.
  \end{itemize}
\item \textbf{Persistence} against sync~1, without its ties: $\Corr(\varphi_k,\varphi_0)$ in 20 bins of $a=\fracpart(|q_k-q_0|)$, each compared with its mean of $1-6a(1-a)$, over bins with at least $10^4$ weights; and the same along all 436 syncs, from the row samples.
\item \textbf{Partial sums,} on the row samples along the recorded path: RTN, stochastic rounding with fresh seeded dither, and error feedback, in units of each row's $\Delta$ at sync~1.
\item \textbf{Policy space.} 128 GSM8K training prompts (a fixed random subset), with one response each sampled at temperature 1 from the initial weights, at most 256 tokens long.
  \begin{itemize}[leftmargin=*,itemsep=1pt]
  \item For $\log\mu$, each linear weight is replaced by its INT8 codes times its row scale. The input of each linear layer is quantized per token (absmax over 127, round half to even), and embeddings, normalization and the output head stay unquantized, as in the rollout engine.
  \item The slope is an OLS fit of $\Delta\log\mu$ on $\Delta\log\pi$ over tokens. Its standard error, 0.001--0.004 for the cumulative pairs, treats the tokens as independent and is therefore too small; no rule uses it.
  \end{itemize}
\item \textbf{Simulation} (exploratory, on a CPU). It uses 22.0M synthetic Gaussian weights in rows of 896, with row scales log-uniform over three octaves.
  \begin{itemize}[leftmargin=*,itemsep=1pt]
  \item Each row's $\nu$ is drawn log-normally to match the 0.5B checkpoint (median $1.78\times10^{-3}$).
  \item The rows are rounded to BF16 or kept in FP32, then moved by Adam on pure-noise gradients, four updates per sync as in the run. This gives predicted flip shares 6--15\% above the run's.
  \item The statistics are those of the frozen analysis script, imported unchanged.
  \end{itemize}
\end{itemize}
